\section{Introduzione}
\subsection{Il panorama della fisica subatomica odierna}
Il \textbf{modello standard} (MS) rappresenta uno dei più grandi raggiungimenti della scienza moderna, in particolare nell'ambito della fisica delle interazioni fondamentali ad alte energie. La sua attendibilità è frutto di un processo che nasce nel cuore del XX secolo e che trova la definitiva consacrazione nel 2012, con la rivelazione del bosone di Higgs H. 
$\\$
Tuttavia, vi sono ancora delle tematiche che non trovano risposta all'interno di questo paradigma, come la massa non banale dei neutrini e la presenza di materia oscura in scala cosmologica. In ogni caso, si può concludere come il MS rappresenti un'ottima descrizione fenomenologica dei processi subatomici a scale energetiche dell'ordine del TeV.
$\\$
All'interno di questo panorama, l'esperimento \textbf{Belle II}, sito all'acceleratore SuperKEKB di Tsukuba, Giappone, ricopre un ruolo fondamentale per esplorare le suddette tematiche, affrontando degli scenari che si spingono aldilà della fisica standard, entrando, pertanto, nel territorio della cosiddetta \emph{nuova fisica}.
$\\$
Parallelamente agli esperimenti che cercano di magnificare l'aspetto energetico, Belle II si colloca ai limiti d'intensità. L'obiettivo principale, dunque, non è quello di aumentare l'energia nel centro di massa, bensì esasperare la luminosità istantanea a valori fino a $8 \times 10^{35} cm^{-2} s^-{1}$. Questo permette di studiare fenomeni estremamente rari e di far emergere eventuali discrepanze con le predizioni del Modello Standard.

\subsection{Macrostruttura dell'elaborato}